\section{Чистые стратегии и игры в нормальной форме}

\subsection{Определение игры в нормальной форме}

В данном разделе вводятся базовые понятия статических игр с полной информацией, где каждый игрок принимает решение однократно и одновременно с другими игроками.

\begin{definition}[Игра в нормальной форме]
Игра в нормальной форме задается тройкой $G = \langle I, S, u \rangle$, где:
\begin{itemize}
    \item $I$ — множество игроков.
    \item $S_i$ — множество стратегий игрока $i$.
    \item $S = \times_{i \in I} S_i$ — множество всех возможных профилей стратегий в игре.
    \item $S_{-i} = \times_{j \neq i} S_j$ — множество профилей стратегий всех игроков, кроме $i$. Профиль стратегий может быть записан как $s = (s_i, s_{-i}) \in S$.
    \item $u_i \colon S \to \mathbb{R}$ — функция выигрыша (полезности) игрока $i$, ставящая в соответствие каждому профилю стратегий вещественное число.
\end{itemize}
\end{definition}

\begin{example}[Камень-ножницы-бумага]
Рассмотрим классическую игру с множеством игроков $I = \{1, 2\}$ и множествами стратегий $S_1 = S_2 = \{\text{«камень»}, \text{«ножницы»}, \text{«бумага»}\}$. Функция выигрышей задается матрицей, представленной в Таблице \ref{tab:rps}.

\begin{table}[h]
\centering
\renewcommand{\arraystretch}{1.3}
\begin{NiceTabular}{cc|ccc}
\toprule
& & \multicolumn{3}{c}{\textbf{Игрок 2}} \\
& & \textit{«камень»} & \textit{«ножницы»} & \textit{«бумага»} \\
\midrule
\multirow{3}{*}{\textbf{Игрок 1}} 
& \textit{«камень»} & $0, 0$ & $1, -1$ & $-1, 1$ \\
& \textit{«ножницы»} & $-1, 1$ & $0, 0$ & $1, -1$ \\
& \textit{«бумага»} & $1, -1$ & $-1, 1$ & $0, 0$ \\
\bottomrule
\end{NiceTabular}
\caption{Матрица выигрышей в игре «Камень-ножницы-бумага»}
\label{tab:rps}
\end{table}
\end{example}

\begin{example}[Дилемма заключенного]
Два подозреваемых, Петя и Вася, находятся в разных камерах и лишены возможности общения. Множество игроков $I = \{\text{Петя}, \text{Вася}\}$, множества стратегий $S_1 = S_2 = \{\text{«сознаться»}, \text{«молчать»}\}$.

\begin{table}[h]
\centering
\renewcommand{\arraystretch}{1.3}
\begin{NiceTabular}{cc|cc}
\toprule
& & \multicolumn{2}{c}{\textbf{Вася}} \\
& & \textit{«молчать»} & \textit{«сознаться»} \\
\midrule
\multirow{2}{*}{\textbf{Петя}} 
& \textit{«молчать»} & $-1, -1$ & $-10, 0$ \\
& \textit{«сознаться»} & $0, -10$ & $-8, -8$ \\
\bottomrule
\end{NiceTabular}
\caption{Дилемма заключенного}
\end{table}
Независимо от выбора оппонента, каждому игроку строго выгоднее выбрать стратегию \textit{«сознаться»}.
\end{example}

\section{Доминирование стратегий}

\begin{definition}[Сильное и слабое доминирование]
Пусть $G = \langle I, S, u \rangle$ — игра в нормальной форме. Для игрока $i$ стратегия $s_i \in S_i$ \textbf{сильно доминирует} стратегию $s_i' \in S_i$, если для всех возможных стратегий остальных игроков $s_{-i} \in S_{-i}$ выполняется строгое неравенство:
\begin{equation}
    u_i(s_i, s_{-i}) > u_i(s_i', s_{-i}).
\end{equation}

Стратегия $s_i \in S_i$ \textbf{слабо доминирует} стратегию $s_i' \in S_i$, если для всех $s_{-i} \in S_{-i}$ выполняется нестрогое неравенство:
\begin{equation}
    u_i(s_i, s_{-i}) \ge u_i(s_i', s_{-i}),
\end{equation}
и хотя бы для одного профиля $s_{-i} \in S_{-i}$ выполняется строгое неравенство $u_i(s_i, s_{-i}) > u_i(s_i', s_{-i})$.
\end{definition}

\begin{definition}[Равновесие в доминирующих стратегиях]
Набор стратегий $s^* \in S$ является равновесием в сильно (слабо) доминирующих стратегиях, если для каждого игрока $i \in I$ и всех его альтернативных стратегий $s_i' \in S_i$ ($s_i' \neq s_i^*$) стратегия $s_i^*$ сильно (слабо) доминирует $s_i'$.
\end{definition}

\begin{example}[Аукцион второй цены]
Продается картина. Два покупателя оценивают ее в $v_1 \in [0, 1]$ и $v_2 \in [0, 1]$ соответственно. Каждый делает заявку $s_i \in [0, 1]$. Товар достается предложившему наибольшую цену, но платит он сумму, равную заявке проигравшего. При совпадении заявок победитель определяется случайным образом с вероятностью $\frac{1}{2}$. 
Функция выигрыша первого игрока:
\begin{equation}
    u_1(s_1, s_2) =
    \begin{cases}
        v_1 - s_2, & s_1 > s_2 \\
        \frac{v_1 - s_1}{2}, & s_1 = s_2 \\
        0, & s_1 < s_2
    \end{cases}
\end{equation}
В данной игре стратегия правдивой ставки $s_i = v_i$ слабо доминирует любую другую стратегию $s_i'$. Например, при $s_2 < v_1$, отклонение к $s_1' < s_2$ приносит $0$ вместо $v_1 - s_2 > 0$. При $s_2 > v_1$, отклонение к $s_1' > s_2$ приносит отрицательный выигрыш $v_1 - s_2 < 0$ вместо $0$.
\end{example}

\begin{example}[Голосование за 2 альтернативы]
Пусть $N$ (нечетное число) собственников выбирают управляющую компанию: «1» (ГБУ Жилищник) или «2» (ТСЖ). Выигрыш собственника $i$ равен $1$, если выбрана компания «1». Стратегия $s_i = 1$ слабо доминирует $s_i' = 2$. Голос игрока имеет значение только тогда, когда голоса остальных разделены поровну (он является решающим). Таким образом:
\begin{equation}
    u_i(1, D_{-i}) - u_i(2, D_{-i}) = 
    \begin{cases}
        1, & \text{если } D_{-i} = \frac{N-1}{2} \\
        0, & \text{иначе}
    \end{cases}
\end{equation}
где $D_{-i}$ — число голосов остальных за первую альтернативу. Следовательно, искреннее (простое) голосование рационально.
\end{example}

\subsection{Итеративное удаление доминируемых стратегий}

Рациональные игроки не будут использовать строго доминируемые стратегии. Предполагая общность знания рациональности, можно последовательно исключать такие стратегии из рассмотрения.

\begin{definition}[Разрешимость по доминированию]
Пусть $G^1, G^2, \dots$ — последовательность игр, полученных путем итеративного удаления сильно доминируемых стратегий. Обозначим $S_i^n$ как множество стратегий игрока $i$ в игре $G^n$, полученное удалением доминируемых стратегий из $S_i^{n-1}$. Пусть $S_i^\infty$ — предел последовательности вложенных множеств $S_i^1 \supseteq S_i^2 \supseteq \dots$. 

Игра называется \textbf{разрешимой по доминированию}, если $S_i^\infty$ состоит ровно из одной стратегии $s_i^\infty$ для каждого $i \in I$. Профиль $s^\infty = (s_1^\infty, \dots, s_N^\infty)$ называется решением игры по доминированию.
\end{definition}

\begin{remark}
Результат итеративного удаления \textbf{слабо} доминируемых стратегий, в отличие от строго доминируемых, в общем случае зависит от порядка их удаления. Использование данного метода требует осторожности.
\end{remark}

\section{Равновесие Нэша}

Во многих играх отсутствуют доминирующие стратегии, однако существуют логически обоснованные исходы.

\begin{example}[Семейный спор]
Муж и жена выбирают, куда пойти: на балет или на бокс.

\begin{table}[h]
\centering
\renewcommand{\arraystretch}{1.3}
\begin{NiceTabular}{cc|cc}
\toprule
& & \multicolumn{2}{c}{\textbf{Жена}} \\
& & \textit{балет} & \textit{бокс} \\
\midrule
\multirow{2}{*}{\textbf{Муж}} 
& \textit{балет} & $3, 5$ & $0, 0$ \\
& \textit{бокс} & $0, 0$ & $4, 2$ \\
\bottomrule
\end{NiceTabular}
\end{table}
Доминирующих стратегий нет. Однако профили (\textit{балет}, \textit{балет}) и (\textit{бокс}, \textit{бокс}) являются взаимоприемлемыми: если один из игроков уверен в выборе другого, у него нет стимула отклоняться.
\end{example}

\begin{definition}[Равновесие Нэша]
Пусть $G = \langle I, S, u \rangle$ — игра в нормальной форме. Профиль стратегий $s^* \in S$ называется равновесием Нэша, если для каждого игрока $i \in I$ и для любой стратегии $s_i' \in S_i$ выполняется:
\begin{equation}
    u_i(s_i^*, s_{-i}^*) \ge u_i(s_i', s_{-i}^*).
\end{equation}
\end{definition}

\begin{remark}[Цитата Р. Майерсона]
«Когда меня спрашивают, почему в какой-то игре игроки должны вести себя как в равновесии Нэша, я отвечаю: "Почему нет?". Далее я предлагаю задавшему вопрос сформулировать свое представление о том, как должны повести себя игроки. Если эта спецификация не является равновесием Нэша, то... она будет противоречить сама себе, если мы предположим, что игроки имеют верное представление о действиях друг друга.»
\end{remark}

\subsection{Функции реакции}

\begin{definition}[Функция наилучшего ответа]
Функция реакции (наилучшего ответа) игрока $i$ есть точечно-множественное отображение $\BR_i \colon S_{-i} \to 2^{S_i}$, определяемое как:
\begin{equation}
    \BR_i(s_{-i}) = \left\{ s_i \in S_i \;\middle|\; u_i(s_i, s_{-i}) = \max_{s_i' \in S_i} u_i(s_i', s_{-i}) \right\}.
\end{equation}
\end{definition}

\begin{lemma}
Профиль стратегий $s^*$ является равновесием Нэша тогда и только тогда, когда для всех $i \in I$ выполняется $s_i^* \in \BR_i(s_{-i}^*)$.
\end{lemma}

\subsection{Доминирование и равновесие Нэша}

\begin{lemma}\label{lem:nash_survives}
Пусть в игре $G = \langle I, S, u \rangle$ профиль $s^* \in S$ является равновесием Нэша. Тогда $s^* \in S^\infty$ (равновесие Нэша всегда выживает при итеративном удалении строго доминируемых стратегий).
\end{lemma}

\begin{proof}
Предположим противное: пусть $s^*$ — равновесие Нэша, не содержащееся в $S^\infty$. Тогда существует первый шаг $n$, на котором стратегия из $s^*$ удаляется. Пусть это стратегия $s_i^* \in S_i^n$. Это означает, что существует стратегия $s_i' \in S_i^n$, которая строго доминирует $s_i^*$ на множестве $S_{-i}^n$:
\begin{equation}\label{eq:dom_proof}
    u_i(s_i', s_{-i}) > u_i(s_i^*, s_{-i}) \quad \forall s_{-i} \in S_{-i}^n.
\end{equation}
Поскольку шаг $n$ — первый, на котором удаляется часть равновесного профиля, профиль стратегий остальных игроков $s_{-i}^*$ все еще принадлежит $S_{-i}^n$. Подставим $s_{-i}^*$ в неравенство \eqref{eq:dom_proof}:
\begin{equation}
    u_i(s_i', s_{-i}^*) > u_i(s_i^*, s_{-i}^*).
\end{equation}
Однако это прямо противоречит определению равновесия Нэша. Следовательно, $s^*$ не может быть удалено.
\end{proof}

\begin{lemma}
Пусть $G = \langle I, S, u \rangle$ — игра с конечным множеством профилей стратегий, и $s^\infty$ — решение этой игры по доминированию (игра разрешима). Тогда $s^\infty$ — единственное равновесие Нэша.
\end{lemma}

\begin{proof}
Предположим противное: множество выживших стратегий $S^\infty = \{s'\}$ не является равновесием Нэша. Тогда существует игрок $i$ и стратегия $s_i \in S_i$, для которых:
\begin{equation}\label{eq:unique_nash_1}
    u_i(s_i, s_{-i}') > u_i(s_i', s_{-i}').
\end{equation}
Так как $s_i$ была удалена, существует шаг $n$ и стратегия $s_i'' \in S_i^n$, которая строго доминирует $s_i$ на $S_{-i}^n$:
\begin{equation}
    u_i(s_i'', s_{-i}) > u_i(s_i, s_{-i}) \quad \forall s_{-i} \in S_{-i}^n.
\end{equation}
Поскольку $s_{-i}' \in S_{-i}^\infty \subseteq S_{-i}^n$, подставим $s_{-i}'$:
\begin{equation}\label{eq:unique_nash_2}
    u_i(s_i'', s_{-i}') > u_i(s_i, s_{-i}').
\end{equation}
Из \eqref{eq:unique_nash_1} и \eqref{eq:unique_nash_2} следует, что $s_i'' \neq s_i'$. В силу конечности множества стратегий, повторяя эти рассуждения для стратегии $s_i''$ (которая также должна быть удалена на некотором шаге $m > n$), мы построим бесконечную цепочку строго доминирующих стратегий, что невозможно в конечной игре. Полученное противоречие доказывает лемму.
\end{proof}

\section{Координационные игры и приложения}

\begin{example}[Списывание на экзамене]
В аудитории $N$ студентов. Множество стратегий $S_i = \{0, 1\}$, где $1$ — «списать», $0$ — «не списывать». Тяжесть наказания обратно пропорциональна числу списывающих (преподаватель не может строго наказать всю группу). Функция выигрыша:
\begin{equation}
    u_i(s) = 
    \begin{cases}
        1 - \frac{C}{\sum_{j=1}^N s_j}, & s_i = 1 \\
        0, & s_i = 0
    \end{cases}
\end{equation}
где $C > 0$ — суммарный объем наказания, $1$ — базовая полезность от списывания. Функция реакции принимает вид:
\begin{equation}
    \BR_i(s_{-i}) =
    \begin{cases}
        0, & \sum_{j \neq i} s_j < C - 1 \\
        \{0, 1\}, & \sum_{j \neq i} s_j = C - 1 \\
        1, & \sum_{j \neq i} s_j > C - 1
    \end{cases}
\end{equation}
Анализ равновесий:
\begin{itemize}
    \item При $C < 1$: единственное равновесие — все списывают.
    \item При $1 \le C \le N$: два чистых равновесия — либо никто не списывает, либо списывают все (возникает самосбывающийся прогноз, зависящий от ожиданий действий других).
    \item При $C > N$: единственное равновесие — никто не списывает.
\end{itemize}
\end{example}

\begin{example}[Конкурс красоты / Угадай $\frac{2}{3}$ среднего]
В игре участвуют $N$ игроков. Стратегическое множество $S_i = [0, 100]$. Выигрыш определяется как:
\begin{equation}
    u_i = - \left| s_i - x \sum_{j \neq i} s_j \right|
\end{equation}
где $x \in (0, 1]$ — параметр. 
При $x = 1$ игра является координационной: любой профиль $s_1 = \dots = s_N = s$ является равновесием.
При $x < 1$ стратегия $s_i = 100x$ доминирует любую стратегию $s_i' \in (100x, 100]$. После первой итерации исключения пространство стратегий сужается до $S_i^1 = [0, 100x]$. На второй итерации — до $S_i^2 = [0, 100x^2]$, и так далее. Решение по доминированию: $s_1 = \dots = s_N = 0$. 

Экономические эксперименты показывают, что люди редко играют $0$, что объясняется ограниченной рациональностью и ожиданием нерационального поведения других участников (слишком большое число шагов итеративного удаления для достижения равновесия).
\end{example}

\begin{example}[Дуополия Курно]
Две фирмы выбирают объемы выпуска $q_i \in [0, \infty)$. Обратная функция спроса: $p = 1 - q_1 - q_2$. Предельные издержки постоянны и равны $c$.
Функции выигрыша (прибыли):
\begin{equation}
    u_i = q_i(1 - q_1 - q_2) - c q_i \quad \text{для } i=1, 2.
\end{equation}
Функция реакции первой фирмы выводится из условия первого порядка:
\begin{equation}
    \BR_1(q_2) = 
    \begin{cases}
        \frac{1 - q_2 - c}{2}, & q_2 < 1 - c \\
        0, & q_2 \ge 1 - c
    \end{cases}
\end{equation}
Процесс итеративного удаления доминируемых стратегий:
\begin{itemize}
    \item Максимально возможный наилучший ответ первой фирмы достигается при $q_2 = 0$ и равен $q_1' = \frac{1-c}{2}$. Все $q_1 > \frac{1-c}{2}$ строго доминируются. Множества выживших стратегий: $S_1^1 = S_2^1 = \left[0, \frac{1-c}{2}\right]$.
    \item Зная, что $q_2 \le \frac{1-c}{2}$, первая фирма понимает, что ее оптимальный выпуск не опустится ниже $q_1'' = \BR_1\left(\frac{1-c}{2}\right) = \frac{1-c}{4}$. Множества стратегий сужаются до $S_1^2 = S_2^2 = \left[\frac{1-c}{4}, \frac{1-c}{2}\right]$.
    \item На следующем шаге верхняя граница сдвигается: $S_1^3 = S_2^3 = \left[\frac{1-c}{4}, \frac{3(1-c)}{8}\right]$.
\end{itemize}
Последовательность интервалов сходится к единственной точке, являющейся равновесием Нэша: $q_1^* = q_2^* = \frac{1-c}{3}$.
\end{example}
